\subsection{THE RIGHT AND LEFT UNIFORMITIES ON A TOPOLOGICAL GROUP}

In a topological group G we can perceive the possibility of defining a notion of ``sufficiently near points'' and hence a uniform structure, by operating as follows: if x and y are any two points of G, we apply to both points the translation which sends one of them, say x, to the identity element e; the ``proximity'' of x and y is then evaluated in some sense by the neighbourhood V of e into which y is translated. This translation, which consists in multiplying both x and y by \(x^{-1}\) , can be carried out on the right or on the left, and we shall see that in either case we obtain effectively a uniformity on G which is compatible with the topology of G. Let us take the case in which the translations are performed on the right; then to each neighbourhood V of e there corresponds the set \(V_{d}\) of pairs \((x, y) \in \mathbf{G} \times \mathbf{G}\) such that \(yx^{-1} \in V\) . Let \(G_{d}\) be the family of sets \(V_{d}\) , as V runs through the neighbourhood filter B of e. Then \(G_{d}\) is a fundamental system of entourages (Chapter II, § I, no. I). For since \(e \in V\) , the diagonal \(\Delta\) of \(G \times G\) is contained in \(V_{d}\) for each \(V \in B\) , hence \(G_{d}\) is a filter base and satisfies axiom \((\mathrm{U}_{\mathrm{I}}^{\prime})\) ; since the relations \(yx^{-1} \in V\) and \(xy^{-1} \in V^{-1}\) are equivalent, we have \(\mathbf{V}_{d}^{-1} = (\mathbf{V}^{-1})_{d}\) .

hence \(\mathbf{V}_{d}^{-1}\in \mathfrak{G}_d\) by \((\mathrm{GV}_{\mathrm{II}})\) , so that \((\mathrm{U}_{\mathrm{II}}^{\prime})\) is satisfied; and finally, the relations \(zx^{-1}\in \mathrm{V}\) and \(yz^{-1}\in \mathrm{V}\) imply \(yx^{-1}\in \mathrm{V.V}\) ; hence \(\mathbf{V}_d\circ \mathbf{V}_d\) is contained in \((\mathrm{V.V})_d\) , and \((\mathrm{GV}_{\mathrm{I}})\) shows that \(\mathfrak{G}_d\) satisfies \((\mathrm{U}_{\mathrm{III}}^{\prime})\) .

The uniformity defined by \(\mathfrak{G}_d\) is compatible with the topology of \(\mathbf{G}\) , for the relations \(y \in \mathrm{V}_d(x)\) and \(y \in \mathrm{V}.x\) are equivalent by definition; in other words \(\mathrm{V}_d(x) = \mathrm{V}.x\) .

The argument is analogous when the translations are on the left, and we may therefore make the following definition:

DEFINITION 1. The right (resp. left) uniformity on a topological group G is the uniformity for which a fundamental system of entourages is obtained by making correspond to each neighbourhood V of the identity element e, the set \(V_{d}\) (resp. \(V_{s}\) ) of pairs \((x, y)\) such that \(y x^{-1} \in V\) (resp. \(x^{-1} y \in V\) ).

If V runs through a fundamental system of neighbourhoods of e, the sets \(V_{d}\) (resp. \(V_{s}\) ) form a fundamental system of entourages of the right (resp. left) uniformity.

To each proposition on the topology of a uniform space there corresponds a proposition on the topology of a group; the translation is made according to Definition 1 and the formulae \(\mathrm{V}_{d}(x)=\mathrm{V}.x,\quad\mathrm{V}_{d}(\mathrm{A})=\mathrm{V}.\mathrm{A},\quad\mathrm{V}_{s}(x)=x.\mathrm{V},\quad\mathrm{V}_{s}(\mathrm{A})=\mathrm{A}.\mathrm{V}\) , which are immediate consequences of the definition. For example, if A is any non-empty subset of G then we have (Chapter II, § 1, no. 2, Corollary 1 to Proposition 2)

\[
\overline {{{\mathrm{A}}}} = \bigcap_ {\mathbf {V} \in \mathfrak {B}} \mathrm{V}. \mathrm{A} = \bigcap_ {\mathbf {V} \in \mathfrak {B}} \mathrm{A}. \mathrm{V}.\tag{1}
\]

Again (Chapter II, § 1, no. 2, Corollary 3 to Proposition 2), every Hausdorff group is regular.

The right and left uniformities on a topological group are in general distinct (see Exercise 4). Obviously they coincide if the group is commutative, for then \(\mathbf{V}_d = \mathbf{V}_s\) ; also they coincide if the group is compact (Chapter II, § 4, no. 1, Theorem 1).

In general, we shall denote by \(G_{s}\) (resp. \(G_{d}\) ) the uniform space obtained by giving the set G its left (resp. right) uniformity.

PROPOSITION 1. The left and right translations are isomorphisms of the right uniformity onto itself.

As to the right translations, the result is clear, since the relation \(yx^{-1} \in V\) is equivalent to \((ya)(xa)^{-1} \in V\) {[}in other words, the mapping \((x, y) \to (xa, ya)\) leaves \(V_d\) fixed{]}. For the left translations, the result follows from \((\mathrm{GV}_{\mathrm{III}})\) ; for \(yx^{-1} \in V\) if and only if \((ay)(ax)^{-1} \in aVa^{-1}\) ; hence \(x \to ax\) is uniformly continuous on \(G_{d}\) .

Similarly, the right and left translations are isomorphisms of the left uniformity onto itself.

Every inner automorphism \(x \rightarrow axa^{-1}\) of G is therefore an automorphism for the group structure of G, for the topology of G, and for both the uniformities of G.

PROPOSITION 2. The symmetry \(x \to x^{-1}\) is an isomorphism of the right uniformity onto the left uniformity.

This is an immediate consequence of Definition 1.

\pandocbounded{\includegraphics[keepaspectratio]{images/part002_77aeac5d.jpg}}

The reader should beware of supposing that the mapping \((x, y) \to xy\) of the uniform space \(G_{d} \times G_{d}\) into the uniform space \(G_{d}\) is in general uniformly continuous. Similarly, the symmetry \(x \to x^{-1}\) , considered as a mapping of \(G_{d}\) onto \(G_{d}\) , is not in general uniformly continuous (see Exercises 3 and 4).

PROPOSITION 3. Every continuous homomorphism \(f\) of a topological group \(G\) into a topological group \(G'\) is uniformly continuous when considered as a mapping of \(G_d\) into \(G_d'\) (or of \(G_s\) into \(G_s'\) ).

For if \(V'\) is a neighbourhood of the identity element in \(G'\) , and \(\mathbf{V} = f^{-1}(\mathbf{V}')\) , then the relation \(yx^{-1} \in V\) implies

\[
f (y) (f (x)) ^ {- 1} = f (y x ^ {- 1}) \in V ^ {\prime}.
\]

